\section{Conclusão e Trabalhos Futuros}
\label{sec:conclusao}

Este trabalho apresentou uma abordagem algorítmica e topológica baseada em grafos heterogêneos multicamada para a identificação de indícios de sobreposição societária, administrativa, cadastral e onomástica entre empresas que disputaram contratações públicas municipais. Aplicada aos dados do SICOM/TCE-MG referentes ao exercício de 2025 de Divinópolis/MG e complementada por informações do QSA obtidas via BrasilAPI em novembro de 2026, a metodologia demonstrou como relações licitatórias e vínculos empresariais podem ser integrados em uma estrutura única e submetidos a critérios determinísticos de seleção.

\subsection{Eficiência da Redução Topológica}

A filtragem por escore multidimensional, com limiar $\theta\geq2{,}00$, seguida da remoção de vértices isolados e da identificação de componentes conexas, reduziu a rede original de 1.788 vértices e 5.914 arestas para um subgrafo com 26 vértices e 23 arestas. A redução corresponde a 98,55\% dos vértices e 99,61\% das arestas. O subgrafo resultante contém seis componentes conexas, formadas por 12 empresas e 14 integrantes do QSA ou titulares de MEI.

Esse resultado demonstra a capacidade da filtragem determinística de concentrar a inspeção em relações selecionadas pelo escore e em seus vínculos incidentes na camada \texttt{É\_Dono\_De}. A redução, entretanto, deve ser entendida como um mecanismo de priorização: ela preserva as relações que satisfazem os critérios definidos, mas não garante que todos os possíveis padrões de risco estejam contidos no subgrafo.

\subsection{Identificação de Sobreposições no QSA e Cadastrais}

A avaliação de 4.059 pares de empresas co-disputantes identificou 202 pares com escore positivo. Desses, 159 enquadraram-se na faixa baixa, 37 na faixa média e seis na faixa alta do escore heurístico de triagem. A camada \texttt{Vinculo\_Suspeito} reteve os 43 pares das duas últimas faixas, enquanto a redução topológica selecionou os seis pares com escore igual ou superior a 2,00.

A enumeração dos motivos triangulares identificou três ocorrências de $K_3$ em dois componentes. Esses motivos representam pares de empresas co-disputantes que compartilham um ou mais integrantes do QSA (sócios ou administradores) na representação agregada multicamada. Os quatro componentes restantes não apresentaram vínculos compartilhados no QSA e foram selecionados por coincidências de endereço, telefone ou, exclusivamente, sobrenomes (Casos 5 e 6). Tais relações são sinais complementares de triagem e exigem validação independente.

Quanto à métrica de recorrência de co-disputa, verificou-se no código que \texttt{Peso\_Concorrencia} soma os certames compartilhados e os itens disputados em comum no julgamento, explicando integralmente o valor 60 no Caso 2 (1 certame + 59 itens) e o valor 65 no Caso 3 (1 certame + 64 itens), bem como o valor 1 nos demais casos pareados exclusivamente no nível do certame.

\subsection{Contribuição ao Controle Externo e Limitações}

O pipeline~\cite{projeto_github} oferece um procedimento rastreável para organizar relações entre órgãos, fornecedores e integrantes do QSA, calcular escores explícitos e reduzir a rede a componentes que podem ser examinadas individualmente. Sua contribuição está na integração de camadas relacionais e na explicabilidade dos critérios de seleção, não na classificação automática de empresas como fraudulentas.

Os resultados possuem caráter indiciário. A presença de vínculos societários, administrativos, cadastrais ou onomásticos não comprova, isoladamente, conluio ou irregularidade. A confirmação de qualquer hipótese exige diligências adicionais, incluindo análise das propostas, dos lances, das atas de julgamento, da execução contratual e de outros elementos documentais pertinentes.

Quatro limitações principais delimitam esta versão: (i) ausência de resultados financeiros validados associados individualmente aos pares, em razão de chaves intermódulo preenchidas com sentinela \texttt{-1} no SICOM; (ii) defasagem temporal do QSA, coletado em 14 de novembro de 2026 em relação a certames ocorridos em 2025; (iii) ausência de filtragem das qualificações do QSA, fazendo com que administradores e diretores não sócios também componham a camada \texttt{É\_Dono\_De}; e (iv) natureza heurística dos pesos e limiares de triagem, ainda sem calibração empírica por casos rotulados ou análise de sensibilidade, bem como possibilidade de coincidências legítimas em endereço, telefone e sobrenomes.

\subsection{Trabalhos Futuros}

Como continuidade, pretende-se desenvolver uma integração validada entre os registros licitatórios e a execução financeira, reconstruindo as relações entre homologações, contratos, empenhos, liquidações e pagamentos. Pretende-se também restringir a formação de pares \texttt{Concorreu\_Contra} ao nível de lote/item, segmentar formalmente as arestas do QSA por qualificação (propriedade versus gestão/administração) e validar a temporalidade histórica dos quadros societários.

Uma direção complementar consiste em incorporar atributos temporais e financeiros das propostas e dos lances registrados nas atas de julgamento, de modo a investigar padrões de alternância de vencedores (\textit{bid rotation}), propostas de cobertura (\textit{cover bidding}) e supressão de propostas. Por fim, pretende-se investigar modelos de aprendizado de representação em grafos heterogêneos (\textit{Heterogeneous Graph Neural Networks}) para calibrar os pesos do escore multidimensional a partir de casos históricos analisados ou julgados pelos órgãos de controle, preservando interpretabilidade e rastreabilidade.
